\documentclass[10pt,a4paper]{amsart}
\usepackage{amsmath,amssymb,amsthm,graphicx}
\usepackage{verbatim}
\newenvironment{reference}{\comment}{\endcomment}
%\newenvironment{reference}{}{}
\newenvironment{slogan}{\comment}{\endcomment}
\newenvironment{history}{\comment}{\endcomment}

% For commutative diagrams we use Xy-pic
\usepackage[all]{xy}

% We use 2cell for 2-commutative diagrams.
\xyoption{2cell}
% Theorem environments.
%
\theoremstyle{plain}
\newtheorem{theorem}[subsection]{Theorem}
\newtheorem{proposition}[subsection]{Proposition}
\newtheorem{lemma}[subsection]{Lemma}

\theoremstyle{definition}
\newtheorem{definition}[subsection]{Definition}
\newtheorem{example}[subsection]{Example}
\newtheorem{exercise}[subsection]{Exercise}
\newtheorem{situation}[subsection]{Situation}

\theoremstyle{remark}
\newtheorem{remark}[subsection]{Remark}
\newtheorem{remarks}[subsection]{Remarks}

\numberwithin{equation}{subsection}

% Macros
%
\def\lim{\mathop{\mathrm{lim}}\nolimits}
\def\colim{\mathop{\mathrm{colim}}\nolimits}
\def\Spec{\mathop{\mathrm{Spec}}}
\def\Hom{\mathop{\mathrm{Hom}}\nolimits}
\def\Ext{\mathop{\mathrm{Ext}}\nolimits}
\def\SheafHom{\mathop{\mathcal{H}\!\mathit{om}}\nolimits}
\def\SheafExt{\mathop{\mathcal{E}\!\mathit{xt}}\nolimits}
\def\Sch{\mathit{Sch}}
\def\Mor{\mathop{\mathrm{Mor}}\nolimits}
\def\Ob{\mathop{\mathrm{Ob}}\nolimits}
\def\Sh{\mathop{\mathit{Sh}}\nolimits}
\def\NL{\mathop{N\!L}\nolimits}
\def\CH{\mathop{\mathrm{CH}}\nolimits}
\def\proetale{{pro\text{-}\acute{e}tale}}
\def\etale{{\acute{e}tale}}
\def\QCoh{\mathit{QCoh}}
\def\Ker{\mathop{\mathrm{Ker}}}
\def\Im{\mathop{\mathrm{Im}}}
\def\Coker{\mathop{\mathrm{Coker}}}
\def\Coim{\mathop{\mathrm{Coim}}}

% Boxtimes
%
\DeclareMathSymbol{\boxtimes}{\mathbin}{AMSa}{"02}

%
% Macros for moduli stacks/spaces
%
\def\QCohstack{\mathcal{QC}\!\mathit{oh}}
\def\Cohstack{\mathcal{C}\!\mathit{oh}}
\def\Spacesstack{\mathcal{S}\!\mathit{paces}}
\def\Quotfunctor{\mathrm{Quot}}
\def\Hilbfunctor{\mathrm{Hilb}}
\def\Curvesstack{\mathcal{C}\!\mathit{urves}}
\def\Polarizedstack{\mathcal{P}\!\mathit{olarized}}
\def\Complexesstack{\mathcal{C}\!\mathit{omplexes}}
% \Pic is the operator that assigns to X its picard group, usage \Pic(X)
% \Picardstack_{X/B} denotes the Picard stack of X over B
% \Picardfunctor_{X/B} denotes the Picard functor of X over B
\def\Pic{\mathop{\mathrm{Pic}}\nolimits}
\def\Picardstack{\mathcal{P}\!\mathit{ic}}
\def\Picardfunctor{\mathrm{Pic}}
\def\Deformationcategory{\mathcal{D}\!\mathit{ef}}
\usepackage{cleveref}
\crefname{originalremark}{original remark}{original remarks}
\Crefname{originalremark}{Original remark}{Original remarks}
\makeatletter
\newlabel{formal-spaces-remark-base-change-variant-Noetherian}{{0GBF}{source}{}{sourcepointer-1}{}}
\newlabel{formal-spaces-remark-base-change-variant-Noetherian@cref}{{[originalremark][0][]0GBF}{[1][0][]source}{}{}{}}
\makeatother
\crefname{originalremark}{original remark}{original remarks}
\Crefname{originalremark}{Original remark}{Original remarks}
\makeatletter
\newlabel{formal-spaces-remark-base-change-variant-variant-Noetherian}{{0GBG}{source}{}{sourcepointer-2}{}}
\newlabel{formal-spaces-remark-base-change-variant-variant-Noetherian@cref}{{[originalremark][0][]0GBG}{[1][0][]source}{}{}{}}
\makeatother
\crefname{originalreference}{original reference}{original references}
\Crefname{originalreference}{Original reference}{Original references}
\makeatletter
\newlabel{lemma-base-change-flat-continuous}{{0GCA}{source}{}{sourcepointer-3}{}}
\newlabel{lemma-base-change-flat-continuous@cref}{{[originalreference][0][]0GCA}{[1][0][]source}{}{}{}}
\makeatother
\crefname{originalreference}{original reference}{original references}
\Crefname{originalreference}{Original reference}{Original references}
\makeatletter
\newlabel{lemma-finite-type-morphisms}{{0GC0}{source}{}{sourcepointer-4}{}}
\newlabel{lemma-finite-type-morphisms@cref}{{[originalreference][0][]0GC0}{[1][0][]source}{}{}{}}
\makeatother
\crefname{originalreference}{original reference}{original references}
\Crefname{originalreference}{Original reference}{Original references}
\makeatletter
\newlabel{lemma-finite-type-red-morphisms}{{0GC7}{source}{}{sourcepointer-5}{}}
\newlabel{lemma-finite-type-red-morphisms@cref}{{[originalreference][0][]0GC7}{[1][0][]source}{}{}{}}
\makeatother
\crefname{originalreference}{original reference}{original references}
\Crefname{originalreference}{Original reference}{Original references}
\makeatletter
\newlabel{lemma-flat-morphisms}{{0GCD}{source}{}{sourcepointer-6}{}}
\newlabel{lemma-flat-morphisms@cref}{{[originalreference][0][]0GCD}{[1][0][]source}{}{}{}}
\makeatother
\title{Flatness is stable under base change of formal algebraic spaces}
\author{The Stacks Project Authors}
\date{stacks-project:restricted.tex:local-snapshot-20261012; GNU FDL 1.2 or later}
\begin{document}
\maketitle
\noindent\textit{Editorial note.} All delivered mathematical text is the author's own, copied from the pinned Stacks Project chapter file restricted.tex; no statement or proof step was rewritten, reordered or removed.\par\medskip
\noindent\textit{Editorial note.} The counted proof is the author's complete argument: the proof delegates to the continuous-family base-change lemma and transfers the outcome through the property-defines-property lemmas of the chapter\par\medskip
\noindent\textit{Editorial note.} The label lemma-base-change-flat names the counted statement, and the proof environment immediately following it in the source is delivered verbatim as its proof.\par\medskip
\noindent\textit{Editorial note.} The counted statement is carried by the author's own lemma environment in the chapter source.\par\medskip
\noindent\textit{Allowed standard prerequisites (editorial).}\begin{itemize}
\item the two base-change comparison remarks of Formal Spaces for Noetherian formal algebraic spaces (formal-spaces-remark-base-change-variant-variant-Noetherian and formal-spaces-remark-base-change-Noetherian)
\item base-change stability of flatness for continuous families and its locally-finite-type companion (lemma-base-change-flat-continuous of the same chapter, parts (1) and (2))
\item flatness and the two finite-type notions as properties of morphisms read through etale covers (lemma-flat-morphisms, lemma-finite-type-red-morphisms, lemma-finite-type-morphisms of the same chapter)
\end{itemize}\medskip
\ifdefined\hypertarget\hypertarget{sourcepointer-1}{}\fi
\noindent\textit{Source reference.} Original remark 0GBF: sources/stacks-chow/tags.txt, source lines 20370--20370. Prerequisite pointer invoked at proof line(s) 3116; it names a standard fact of a sibling chapter that the counted argument uses, and every step of the counted proof built on it is retained verbatim in the delivered body.\par\medskip
\ifdefined\hypertarget\hypertarget{sourcepointer-2}{}\fi
\noindent\textit{Source reference.} Original remark 0GBG: sources/stacks-chow/tags.txt, source lines 20371--20371. Prerequisite pointer invoked at proof line(s) 3108; it names a standard fact of a sibling chapter that the counted argument uses, and every step of the counted proof built on it is retained verbatim in the delivered body.\par\medskip
\ifdefined\hypertarget\hypertarget{sourcepointer-3}{}\fi
\noindent\textit{Source reference.} Original reference 0GCA: sources/stacks-restricted/pinned-original-source.tex, source lines 2951--2951. Prerequisite pointer invoked at proof line(s) 3109; it names a standard fact of the same chapter that the counted argument uses, and every step of the counted proof built on it is retained verbatim in the delivered body.\par\medskip
\ifdefined\hypertarget\hypertarget{sourcepointer-4}{}\fi
\noindent\textit{Source reference.} Original reference 0GC0: sources/stacks-restricted/pinned-original-source.tex, source lines 2556--2556. Prerequisite pointer invoked at proof line(s) 3119; it names a standard fact of the same chapter that the counted argument uses, and every step of the counted proof built on it is retained verbatim in the delivered body.\par\medskip
\ifdefined\hypertarget\hypertarget{sourcepointer-5}{}\fi
\noindent\textit{Source reference.} Original reference 0GC7: sources/stacks-restricted/pinned-original-source.tex, source lines 2795--2795. Prerequisite pointer invoked at proof line(s) 3111; it names a standard fact of the same chapter that the counted argument uses, and every step of the counted proof built on it is retained verbatim in the delivered body.\par\medskip
\ifdefined\hypertarget\hypertarget{sourcepointer-6}{}\fi
\noindent\textit{Source reference.} Original reference 0GCD: sources/stacks-restricted/pinned-original-source.tex, source lines 3043--3043. Prerequisite pointer invoked at proof line(s) 3110; it names a standard fact of the same chapter that the counted argument uses, and every step of the counted proof built on it is retained verbatim in the delivered body.\par\medskip
\begin{lemma}
\label{lemma-base-change-flat}
Let $S$ be a scheme. Let $f : X \to Y$ and $g : Z \to Y$
be morphisms of locally Noetherian formal algebraic spaces over $S$.
\begin{enumerate}
\item If $f$ is flat and $g_{red} : Z_{red} \to Y_{red}$ is
locally of finite type, then the base change
$X \times_Y Z \to Z$ is flat.
\item If $f$ is flat and locally of finite type, then
the base change $X \times_Y Z \to Z$ is flat and locally of finite type.
\end{enumerate}
\end{lemma}

\begin{proof}
Part (1) follows from a combination of
Formal Spaces, Remark
\ref{formal-spaces-remark-base-change-variant-variant-Noetherian},
Lemma \ref{lemma-base-change-flat-continuous} part (1),
Lemma \ref{lemma-flat-morphisms}, and
Lemma \ref{lemma-finite-type-red-morphisms}.

\medskip\noindent
Part (2) follows from a combination of
Formal Spaces, Remark
\ref{formal-spaces-remark-base-change-variant-Noetherian},
Lemma \ref{lemma-base-change-flat-continuous} part (2),
Lemma \ref{lemma-flat-morphisms}, and
Lemma \ref{lemma-finite-type-morphisms}.
\end{proof}

\end{document}
